\documentclass[UTF8]{ctexart}

\usepackage{amsmath, amssymb, amsthm}
\usepackage{mathrsfs}
\usepackage{ctex}
\usepackage{tikz}
\usepackage{tikz-cd}
\usepackage{pgfplots}
\usepackage{xltabular}
\usepackage{booktabs}
\usepackage{geometry}
\usepackage{extarrows}
\geometry{a4paper, margin=2cm}
\pgfplotsset{compat=1.18}
\usetikzlibrary{shapes.geometric, positioning, calc, decorations.pathreplacing,3d ,positioning}

\newtheorem{theorem}{定理}[subsection]
\newtheorem{proposition}[theorem]{命题}
\newtheorem{corollary}[theorem]{推论}
\newtheorem{lemma}[theorem]{引理}
\theoremstyle{definition}
\newtheorem{definition}[theorem]{定义}
\newtheoremstyle{examplestyle}{3pt}{3pt}{\itshape}{0pt}{\bfseries}{.}{.5em}{\thmnumber{例 \arabic{example}}\thmnote{\ (#3)}}
\theoremstyle{examplestyle}
\newtheorem{example}{}[subsection]
\theoremstyle{remark}
\newtheorem*{remark}{注}

\begin{document}
%###——————————————4.1——————————————————###
原问题等价于存在参数$\lambda>0$, 使得lasso问题
\[min\quad \frac{1}{2}\|Ax-y\|_2^2+\lambda\|x\|_1\]
的最优解同时满足残差约束$\|Ax_\lambda-y\|_2=\delta$, 残差函数$r(\lambda)$不严格单调, 欠定条件下朴素二分法失效, 构造$\lambda$几何递减序列($[10^{-6}\|A^Ty\|_\infty,\|A^Ty\|_\infty]$), 对每个$\lambda$求解lasso的残差曲线, 利用插值寻找$r(\lambda)=\delta$的目标$\lambda$.

%###——————————————4.2——————————————————###
$\|x\|_1$的不可微性, 使得我们无法简单的梯度下降求解, 所以我们引入近端算子, 它是处理不可微凸函数的一种重要工具, 对于一个闭、凸函数
\[f:\mathbb{R}^n\to(-\infty,+\infty]\]
给定参数$\tau>0$, 近端算子定义为
\[\text{prox}_{\tau f}(v)=\arg\min_x\{f(x)+\frac{1}{2\tau}\|x-v\|_2^2\}\]
可以理解成, 在尽量保持$x$接近$v$的同时, 使得$f(x)$尽量小.

我们研究的lasso
\[min\quad \frac{1}{2}\|Ax-y\|_2^2+\lambda\|x\|_1\]
可以看做把目标函数拆成两部分$g(x)=\frac{1}{2}\|Ax-y\|_2^2$, $h(x)=\lambda\|x\|_1$, lasso是二者之和. 其中$g$可微而$h$不可微但凸, 所以我们处理的方向就是梯度下降处理$g$, 近端算子处理$h$, 所以我们要看清楚$\text{prox}_{\tau \lambda\|\cdot\|_1}(v)$是什么.

考虑
\[h(x)=\lambda\|x\|_1\]
根据定义
\[\text{prox}_{\tau h}(v)=\arg\min_x\{\lambda\tau\|x\|_1\}+\frac{1}{2}\|x-v\|_2^2\]
所以我们要求解
\[\min_x\quad\frac{1}{2}\|x-v\|_2^2+\tau\lambda\|x\|_1\]
将目标函数写开来
\[\|x-v\|_2^2=\sum_{i=1}^n(x_i-v_i)\]\[\|x\|_1=\sum_{i=1}^n|x_i|\]
把求和号合并,
\[\sum_{i=1}^n[\frac{1}{2}(x_i-v_i)^2+\tau\lambda|x_i|]\]
所以这个n维优化问题可以分解为n个独立的一维问题, 我们只需研究标量问题
\[\min_x\quad\phi(x)=\frac{1}{2}(x-v)^2+\alpha|x|\]
这里为了简化记号, 令$\alpha=\tau\lambda$.

由于绝对值函数在0处不可微, 分情况讨论.\\
1.$x>0$, 此时$\phi(x)=\frac{1}{2}(x-v)^2+\alpha x$, 直接求导, 容易知道$x^*=v-\alpha$(需要满足$v-\alpha>0$)\\
2.$x<0$, 此时$\phi(x)=\frac{1}{2}(x-v)^2-\alpha x$, 依然直接求导, 在$v<-\alpha$时, $x^*=v+\alpha$.\\
3.x=0, 此时剩下的情况是$-\alpha\le v\le\alpha$, 用次梯度
\[\partial |x|\bigg|_{x=0}=[-1,1]\]
故
\[\partial\phi(0)=-v+\alpha[-1,1]\]
最优性条件是
\[0\in\partial\phi(0)\]
这就是剩下的$-\alpha\le v\le\alpha$, 此时$x^*=0$.

经过分类讨论
\[x^*=\begin{cases}
    v-\alpha, &v>\alpha\\
    0, &|v|\le\alpha\\
    v+\alpha, &v<-\alpha
\end{cases}\]
回代$\alpha=\tau\lambda$, 并把分段函数写成sign格式
\[x^*=\text{sign}(v)\max\{|v|-\tau\lambda,0\}\]
这被称为软阈值算子
\[S_{\tau\lambda}(v)=\text{sign}(v)\max\{|v|-\tau\lambda,0\}\]

现在把标量问题还原到向量, 软阈值算子逐元素作用
\[\text{prox}_{\tau\lambda\|\cdot\|_1}(v)=(S_{\tau\lambda}(v_1),\cdots,S_{\tau\lambda}(v_n))^T\]

现在推导近端梯度法

对于$g(x)$, 在$x_k$点做一阶Taylor近似
\[g(x)=g(x_k)+\nabla g(x_k)^T(x-x_k)\]
于是在$x_k$附近, 我们可以把$g(x)+h(x)$近似成
\[g(x)=g(x_k)+\nabla g(x_k)^T(x-x_k)+h(x)\]
直接最小化它, 它通常没有良好数值稳定性, 于是加上二次正则项$\frac{1}{2t}\|x-x_k\|_2^2$, 另外$g(x_k)$与$x$无关, 在$\arg\min_x$下可以删掉
\[x_{k+1}=\arg\min_x\{\nabla f(x_k)^T(x-x_k)+\frac{1}{2t}\|x-x_k\|_2^2+h(x)\}\]
将$\nabla g(x_k)^T(x-x_k)+\frac{1}{2t}\|x-x_k\|_2^2$配方为
\[\frac{1}{2t}\|x-(x_k-t\nabla g(x_k))\|_2^2+\text{constant}\]
同样地, 在$\arg\min_x$下constant可以去掉, 变为
\[x_{k+1}=\arg\min_x\{h(x)+\frac{1}{2t}\|x-(x_k-t\nabla g(x_k))\|_2^2\}\]
这就变成刚刚定义的近端算子形式, 也就是近端梯度法的迭代公式
\[x_{k+1}=\text{prox}_{th}(x_k-t\nabla g(x_k))\]

现在把$h(x)=\lambda\|x\|_1$代入
\[x_{k+1}=\text{prox}_{t\lambda\|\cdot\|_1}(x_k-t\nabla g(x_k))\]
我们已经证明
\[\text{prox}_{t\|\lambda\|_1}(v)=S_{t\lambda}(v)\]
同时
\[\nabla g(x_k)=A^T(Ax_k-y)\]
所以
\[x_{k+1}=S_{t\lambda}(x_k-tA^T(Ax_k-y))\]
这就是ISTA. 这里t是步长, 取值范围$0<t\le\frac{1}{L}$, $L$是$g$的梯度Lipschitz常数. 原因见下文收敛性.

我们发现ISTA每一步都依靠$x_k$计算, 只利用当前点信息, 现在我们介入Nesterov加速思想, 构造预测点来加速迭代.

构造新的点$y_k$, 使得近端梯度迭代变为
\[x_{k+1}=\text{prox}_{th}(y_k-t\nabla g(y_k))\]
引入辅助变量$t_k$, 定义
\[t_{k+1}=\frac{1+\sqrt{1+4t_k^2}}{2}\]
然后
\[y_{k+1}=x_{k+1}+\frac{t_k-1}{t_{k+1}}(x_{k+1}-x_k)\]
这就是FISTA.

收敛性:

先考虑ISTA. 设$x^*$是最优解, 考察$g$的梯度Lipschitz性:
\[\|\nabla g(x)-\nabla g(z)\|_2=\|A^TA(x-z)\|_2\le\|A^TA\|_2\|x-z\|_2\]
所以具有Lipschitz性, $L=\|A^TA\|_2=\|A\|_2^2$. 步长$t\le\frac{1}{L}=\frac{1}{\|A\|_2^2}$, 由于梯度Lipschitz性
\[\begin{aligned}g(x)&\le g(z)+\nabla g(z)^T(x-z)+\frac{L}{2}\|x-z\|_2^2\\&\le g(z)+\nabla g(z)^T(x-z)+\frac{1}{2t}\|x-z\|^2\end{aligned}\]
定义
\[Q_t(x,z)=g(z)+\nabla g(z)^T(x-z)+\frac{1}{2t}\|x-z\|^2+h(z)\]
就有
\[F(z)\le Q_t(x,z)\]
也就是说$Q_t(x,z)$是$F(x)$的一个上界, 而ISTA的迭代
\[x_{k+1}=\text{prox}_{th}(x_k-t\nabla g(x_k))\]
等价于
\[x_{k+1}=\arg\min_xQ_t(x,x_k)\]
所以
\[0\in\nabla g(x_k)+\frac{1}{t}(x_{k+1}-x_k)+\partial h(x_{k+1})\]
即存在$s_{k+1}\in\partial h(x_{k+1})$使得
\[\nabla g(x_k)+\frac{1}{t}(x_{k+1}-x_k)+s_{k+1}=0\]
即
\[s_{k+1}=-\nabla g(x_k)-\frac{1}{t}(x_{k+1}-x_k)\]
由$h$的凸性,
\[h(x)\ge h(x_{k+1})+s_{k+1}^T(x-x_{k+1})\]
移项并带入$s_{k+1}$
\[h(x_{k+1})-h(x)\le-\nabla g(x_k)^T(x_{k+1}-x)-\frac{1}{t}(x_{k+1}-x_k)^T(x_{k+1}-x)\]
另一方面, 我们一开始就由Lipschitz性得到
\[g(x)\le g(z)+\nabla g(z)^T(x-z)+\frac{1}{2t}\|x-z\|^2\]
我们用$x_{k+1}$替代$x$, 用$x_k$替代$z$, 得到
\[g(x_{k+1})\le g(x_k)+\nabla g(x_k)^T(x_{k+1}-x_k)+\frac{1}{2t}\|x_{k+1}-x_k\|^2\]
结合两个不等式得到
\[F(x_{k+1})-F(x)\le\nabla g(x_k)^T(x_{k+1}-x)+\frac{1}{2t}\|x_{k+1}-x_k\|^2-\frac{1}{t}(x_{k+1}-x_k)^T(x_{k+1}-x)\]
利用恒等式
\[2(a-b)^(a-c)=\|a-b\|^2+\|a-c\|^2-\|b-c\|^2\]
令$a=x_{k+1}$, $b=x_k$, $c=x$得到
\[2(x_{k+1}-x_k)^T(x_{k+1}-x)=\|x_{k+1}-x_k\|^2+\|x_{k+1}-x\|^2-\|x_k-x\|^2\]
代入上面的不等式并整理, 得到
\[F(x_{k+1})-F(x)\le\frac{1}{2t}(\|x_k-x\|^2-\|x_{k+1}-x\|^2)\]
令$x=x^*$, 得到
\[F(x_{k+1})-F^*\le\frac{1}{2t}(\|x_k-x^*\|^2-\|x_{k+1}-x^*\|^2)\]
对$k$从$0$到$N-1$求和
\[\sum_{k=0}^{N-1}(F(x_{k-1})-F^*)\le\frac{1}{2t}\sum_{k=0}^{N-1}(\|x_k-x^*\|^2-\|x_{k+1}-x^*\|^2)\]
右边可以逐项抵消, 只剩
\[\sum_{k=0}^{N-1}(F(x_{k-1})-F^*)\le\frac{1}{2t}\|x_0-x^*\|^2\]
而$F(x_k)$逐项下降
\[F(x_N)-F^*\le\frac{1}{N}\sum_{k=0}^{N-1}(F(x_{k-1})-F^*)\]
所以
\[F(x_N)-F^*\le\frac{\|x_0-x^*\|^2}{2tN}\]
如上, 我们证明了
\[F(x_N)-F^*=O(\frac{1}{N})\]
所以ISTA是$O(\frac{1}{k})$收敛的.


按照我们上述讨论, 步长直接选取$t=\frac{1}{\|A\|_2^2}$, 其中$\|A\|_2^2$通过计算$A^TA$最大特征值求得.

停机标准: 双重准则\\
相对迭代变化$r_x^{(k)}=\frac{\|x_{k+1}-x_k\|_2}{\max(1,\|x_k\|_2)}<\varepsilon_x=10^{-5}$\\
AND 相对目标函数变化$r_F^{(k)}=\frac{|F_{k+1}-F_k|}{\max(1,|F_k|)}<\varepsilon_F=10^{-8}$\\
同时迭代次数$k\ge k_{\text{max}}=500$时强制停机.
\end{document}